\documentclass[10pt,a4paper]{amsart}
\usepackage[margin=19mm]{geometry}
\usepackage{amsmath,amssymb,mathtools}
\usepackage[scr=rsfs,cal=euler]{mathalfa}
\usepackage{xfrac}
\usepackage[unicode,hidelinks,bookmarks=false]{hyperref}
\newcommand{\ie}{i.e.\ }
\newcommand{\cf}{cf.\ }
\newcommand{\resp}{respectively\ }
\newcommand{\Subsection}{\subsection}
\providecommand{\nobreakdash}{\nobreak\hbox{-}}
\setlength{\emergencystretch}{2em}
\multlinegap0pt
\usepackage{bm}
\usepackage{bbm}
\usepackage{dsfont}
\usepackage{xspace}
\usepackage{cancel}
\usepackage{mathtools}
\usepackage{amsthm}
\makeatletter
\@ifundefined{thm}{\newtheorem{thm}{Theorem}}{}
\@ifundefined{lem}{\newtheorem{lem}[thm]{Lemma}}{}
\makeatother
\title[Inverses of six classes of permutation polynomials of the fo]{Inverses of six classes of permutation polynomials of the form $x+\gamma\operatorname{Tr}\_q^{q^2}(h(x))$ over finite fields of even characteristic}
\author{Rajesh P. Singh, Dinesh Kumar, Jitendra Prakash}
\date{Source: arXiv:2601.02919v2 (CC BY 4.0)}
\begin{document}
\maketitle

\noindent\emph{Source and licence.} Inverses of six classes of permutation polynomials of the form $x+\gamma\operatorname{Tr}_q^{q^2}(h(x))$ over finite fields of even characteristic. Version arXiv:2601.02919v2 (CC BY 4.0), distributed under the Creative Commons Attribution 4.0 International licence, as recorded in the arXiv licence metadata for this exact version and shown on the abstract page https://arxiv.org/abs/2601.02919v2. Full rights notice: licenses/arxiv-2601.02919-CC-BY-4.0.txt.\par\smallskip

\noindent\textit{Editorial scope.} Counted unit: the Thm whose source label is \texttt{None} in the article (source0.tex lines 423--429, source label \texttt{None}), delivered together with the article's own complete original proof of it (source0.tex lines 430--588, one \texttt{proof} environment of 159 lines). No mathematics is written, repaired or re-derived: statement and proof are the article's own text line for line. Required context retained uncounted: lem2 (lines 152--157); lem3 (lines 208--212). Every definition, display and label of those lines is retained unchanged. Omitted from this package and not redistributed: 1--151, 158--207, 213--422, 589--951. Nothing inside a retained excerpt is omitted or altered: every retained body is delivered exactly as the article's lines are written, so no omission receipt of any kind accompanies this package. Citations: the retained excerpts carry no citation command at all, so no bibliography entry of the article is transcribed and no reference is redistributed. Reference adaptation: none was needed and none was made; every reference inside the delivered unit resolves inside it, so no unresolved reference or citation is produced. Printed slips kept literal: no printed word, symbol, hypothesis or number of the counted statement or of its proof was altered. Layout and class: the article's own class is \texttt{birkjour} with packages including cite, amsmath, amssymb, amsthm, etoolbox, xcolor, graphicx, hyperref; the wrapper is the lane's pinned \texttt{amsart} editorial wrapper at 10pt on A4, which restates the theorem-like environments the retained text needs and adds the article's own macro definitions that the retained text needs (none), with extra wrapper packages (bm, bbm, dsfont, xspace, cancel, mathtools). The delivered numbering is the unit's own. Assets: no retained span contains a figure environment, a graphics-inclusion command, a PDF-page inclusion, a table or a TikZ picture; no image file is copied into this package, the package declares no asset dependency and no image file is redistributed with it. Licence: version arXiv:2601.02919v2 of this article is distributed under the Creative Commons Attribution 4.0 International licence, as recorded in the arXiv licence metadata for this exact version and shown on the abstract page https://arxiv.org/abs/2601.02919v2. A full rights notice is delivered at licenses/arxiv-2601.02919-CC-BY-4.0.txt. Characters: every retained excerpt is pure ASCII, so the delivered engine, XeLaTeX with its default Latin Modern text font, reports no missing character.
	\begin{lem}\label{lem2}
		Let $q=2^m$ and $\alpha\in\mathbb{F}_{q^2}\setminus\mathbb{F}_q$ with $\operatorname{Tr}_q^{q^2}(\alpha)=1$.  Let $\beta, \delta \in \mathbb{F}^{*}_{q^2}$ be such that $\beta \neq c\delta$ for any $c \in \mathbb{F}^{*}_{q}$. 
		Then 
		$T(x)=\operatorname{Tr}_{q}^{q^2}(\beta x)+\alpha\operatorname{Tr}_{q}^{q^2}(\delta x)$ is a permutation polynomial over $\mathbb{F}_{q^2}$ and its inverse is
		$$T^{-1}(x)=(\beta\delta^{q}+\delta\beta^{q})^{-1}\{\delta^{q}\operatorname{Tr}_{q}^{q^2}((1+\alpha)x)+\beta^{q}\operatorname{Tr}_{q}^{q^2}(x)\}.$$
	\end{lem}

	\begin{lem} \label{lem3}
		Let $q=2^m$, $\alpha\in\mathbb{F}_{q^2}\setminus\mathbb{F}_q$ be such that $\operatorname{Tr}_q^{q^2}(\alpha)=1$ . Then, we have 
		$\operatorname{Tr}_{q}^{q^2}(\beta T^{-1}(x))=\operatorname{Tr}_{q}^{q^2}((1+\alpha)x)$ and 
		$\operatorname{Tr}_{q}^{q^2}(\delta T^{-1}(x))=\operatorname{Tr}_{q}^{q^2}(x).$
	\end{lem}

\begin{thm} The compositional inverse of permutation polynomial 
		$f_3(x)=x+\gamma\operatorname{Tr}_{q}^{q^2}(x+x^3+x^{q+2})$ over $\mathbb{F}_{q^2}$ is \\
	$f_3^{-1}(x)=\begin{cases}x+\gamma\operatorname{Tr}_{q}^{q^2}(x)+\gamma(\operatorname{Tr}_{q}^{q^2}(x))^3,& \text{if}~\gamma\in\mathbb{F}_q\\
	\operatorname{Tr}_{q}^{q^2}(1+\gamma)x+\gamma(\operatorname{Tr}_{q}^{q^2}(x))^t,& \text{if}~\operatorname{Tr}_{q}^{q^2}(\gamma)=1\text{, } \text{and}~ m \text{ is odd}
	\end{cases}$\\
	where $3t\equiv 1\mod{(2^m-1)}$.
\end{thm}

\begin{proof} We can easily see that $\operatorname{Tr}_{q}^{q^2}(\gamma+\gamma^2)=0$ means either $\operatorname{Tr}_{q}^{q^2}(\gamma)=0$ or $\operatorname{Tr}_{q}^{q^2}(\gamma)=1$. First, we assume that $\operatorname{Tr}_{q}^{q^2}(\gamma)=0$ or equivalently $\gamma \in \mathbb{F}_{q}$. We have,
	\begin{align}
		f_3\circ T(x)=&\operatorname{Tr}_{q}^{q^2}(\beta x)+\alpha\operatorname{Tr}_{q}^{q^2}(\delta x)+\gamma\operatorname{Tr}_{q}^{q^2}[\operatorname{Tr}_{q}^{q^2}(\beta x)+\alpha\operatorname{Tr}_{q}^{q^2}(\delta x)+\{\operatorname{Tr}_{q}^{q^2}(\beta x)\notag\\&+\alpha\operatorname{Tr}_{q}^{q^2}(\delta x)\}^3+\{\operatorname{Tr}_{q}^{q^2}(\beta x)+\alpha\operatorname{Tr}_{q}^{q^2}(\delta x)\}^{q+2}]\notag\\
		=&\operatorname{Tr}_{q}^{q^2}(\beta x)+\alpha\operatorname{Tr}_{q}^{q^2}(\delta x)+\gamma\operatorname{Tr}_{q}^{q^2}[\operatorname{Tr}_{q}^{q^2}(\beta x)+\alpha\operatorname{Tr}_{q}^{q^2}(\delta x)+\{\operatorname{Tr}_{q}^{q^2}(\beta x)\notag\\&+\alpha\operatorname{Tr}_{q}^{q^2}(\delta x)\}^2\{\operatorname{Tr}_{q}^{q^2}(\beta x)+\alpha\operatorname{Tr}_{q}^{q^2}(\delta x)+\operatorname{Tr}_{q}^{q^2}(\beta x)+\alpha\operatorname{Tr}_{q}^{q^2}(\delta x)\notag\\&+\operatorname{Tr}_{q}^{q^2}(\delta x)\}]\notag\\
		=&\operatorname{Tr}_{q}^{q^2}(\beta x)+\alpha\operatorname{Tr}_{q}^{q^2}(\delta x)+\gamma\{\operatorname{Tr}_{q}^{q^2}(\delta x)+(\operatorname{Tr}_{q}^{q^2}(\delta x))^3\}.\label{3a}
	\end{align}
	On considering the equation $f_3 \circ T(x)=T(y)=\operatorname{Tr}_{q}^{q^2}(\beta y)+\alpha\operatorname{Tr}_{q}^{q^2}(\delta y)$, we get the following system of equations
	\begin{align*}
		\operatorname{Tr}_{q}^{q^2}(\delta x)=&\operatorname{Tr}_{q}^{q^2}(\delta y)\\
		\operatorname{Tr}_{q}^{q^2}(\beta x)=&\operatorname{Tr}_{q}^{q^2}(\beta y)+\gamma\{\operatorname{Tr}_{q}^{q^2}(\delta y)+(\operatorname{Tr}_{q}^{q^2}(\delta y))^3\}.
	\end{align*}
	Equivalently in matrix form we get,
	$$\begin{bmatrix}
		\beta&\beta^{q}\\
		\delta&\delta^{q}
	\end{bmatrix}
	\begin{bmatrix}
		x\\x^{q}
	\end{bmatrix}=
	\begin{bmatrix}
		\operatorname{Tr}_{q}^{q^2}(\beta y)+\gamma\{\operatorname{Tr}_{q}^{q^2}(\delta y)+(\operatorname{Tr}_{q}^{q^2}(\delta y))^3\}\\
		\operatorname{Tr}_{q}^{q^2}(\delta y)
	\end{bmatrix}.$$
	Applying elementary row operations $R_1\rightarrow\delta^{q}R_1$ and $R_2\rightarrow\beta^{q}R_2$, we get
	$$\begin{bmatrix}
		\beta\delta^{q}&\beta^{q}\delta^{q}\\
		\beta^{q}\delta&\beta^{q}\delta^{q}
	\end{bmatrix}
	\begin{bmatrix}
		x\\x^{q}
	\end{bmatrix}=
	\begin{bmatrix}
		\delta^{q}\operatorname{Tr}_{q}^{q^2}(\beta y)+\delta^{q}\gamma\{\operatorname{Tr}_{q}^{q^2}(\delta y)+(\operatorname{Tr}_{q}^{q^2}(\delta y))^3\}\\
		\beta^{q}\operatorname{Tr}_{q}^{q^2}(\delta y)
	\end{bmatrix}.$$
	Again applying elementary row operations $R_1\rightarrow R_1+R_2$, we get
	\begin{align}\begin{bmatrix}
		\beta\delta^{q}+\beta^{q}\delta&0\\
		\beta^{q}\delta&\beta^{q}\delta^{q}
	\end{bmatrix}
	\begin{bmatrix}
		x\\x^{q}
	\end{bmatrix}=C\label{3b}\end{align}
	\text{where}, $C=
	\begin{bmatrix}
		\delta^{q}\operatorname{Tr}_{q}^{q^2}(\beta y)+\delta^{q}\gamma\{\operatorname{Tr}_{q}^{q^2}(\delta y)+(\operatorname{Tr}_{q}^{q^2}(\delta y))^3\}+\beta^{q}\operatorname{Tr}_{q}^{q^2}(\delta y)\\
		\beta^{q}\operatorname{Tr}_{q}^{q^2}(\delta y)
	\end{bmatrix}.$\\
		From the equation \ref{3b}, we get
		\begin{align*}
	x=&(\beta\delta^{q}+\beta^{q}\delta)^{-1}\{\delta^{q}\operatorname{Tr}_{q}^{q^2}(\beta y)+\delta^{q}\gamma\operatorname{Tr}_{q}^{q^2}(\delta y)+\delta^{q}\gamma(\operatorname{Tr}_{q}^{q^2}(\delta y))^3\\&+\beta^{q}\operatorname{Tr}_{q}^{q^2}(\delta y)\}
\end{align*}
	or
	\begin{align*}
	T^{-1}\circ f_3^{-1}\circ T(y)=&(\beta\delta^{q}+\beta^{q}\delta)^{-1}\{\delta^{q}\operatorname{Tr}_{q}^{q^2}(\beta y)+\delta^{q}\gamma\operatorname{Tr}_{q}^{q^2}(\delta y)+\delta^{q}\gamma\\&(\operatorname{Tr}_{q}^{q^2}(\delta y))^3+\beta^{q}\operatorname{Tr}_{q}^{q^2}(\delta y)\}.
\end{align*}
	We replace $y$ with $T^{-1}(x)$ in above equation
	\begin{align*}
		T^{-1}\circ f_3^{-1}(x)=&(\beta\delta^{q}+\beta^{q}\delta)^{-1}\{\delta^{q}\operatorname{Tr}_{q}^{q^2}(\beta T^{-1}(x))+\delta^{q}\gamma\operatorname{Tr}_{q}^{q^2}(\delta T^{-1}(x))\\&+\delta^{q}\gamma(\operatorname{Tr}_{q}^{q^2}(\delta T^{-1}(x)))^3+\beta^{q}\operatorname{Tr}_{q}^{q^2}(\delta T^{-1}(x))\}.
	\end{align*}
	Now using lemma \ref{lem3} we get
	\begin{align*}
		T^{-1}\circ f_3^{-1}(x)=&(\beta\delta^{q}+\beta^{q}\delta)^{-1}\{\delta^{q}\operatorname{Tr}_{q}^{q^2}((1+\alpha)x)+\delta^{q}\gamma\operatorname{Tr}_{q}^{q^2}(x)\\&+\delta^{q}\gamma(\operatorname{Tr}_{q}^{q^2}(x))^3+\beta^{q}\operatorname{Tr}_{q}^{q^2}(x)\}.
	\end{align*}
	 We know that
	\begin{align}
		f_3^{-1}(x)=&\operatorname{Tr}_{q}^{q^2}[\beta (T^{-1}\circ f_3^{-1}(x))]+\alpha\operatorname{Tr}_{q}^{q^2}[\delta (T^{-1}\circ f_3^{-1}(x))].
		\label{3c}
	\end{align}
		Now we have
	\begin{align}
		\operatorname{Tr}_{q}^{q^2}[\beta (T^{-1}\circ f_3^{-1}(x))]=&\operatorname{Tr}_{q}^{q^2}[\beta (\beta\delta^{q}+\beta^{q}\delta)^{-1}\{\delta^{q}\operatorname{Tr}_{q}^{q^2}((1+\alpha)x)\notag\\&+\delta^{q}\gamma\operatorname{Tr}_{q}^{q^2}(x)+\delta^{q}\gamma(\operatorname{Tr}_{q}^{q^2}(x))^3+\beta^{q}\operatorname{Tr}_{q}^{q^2}(x)\}]\notag\\
		=&\operatorname{Tr}_{q}^{q^2}[\beta\delta^{q}(\beta\delta^{q}+\beta^{q}\delta)^{-1}\operatorname{Tr}_{q}^{q^2}((1+\alpha)x)\notag\\&+\beta\delta^{q}\gamma(\beta\delta^{q}+\beta^{q}\delta)^{-1}\operatorname{Tr}_{q}^{q^2}(x)\notag\\&+\beta\delta^{q}\gamma(\beta\delta^{q}+\beta^{q}\delta)^{-1}(\operatorname{Tr}_{q}^{q^2}(x))^3\notag\\&+\beta\beta^{q}(\beta\delta^{q}+\beta^{q}\delta)^{-1}\operatorname{Tr}_{q}^{q^2}(x)]\notag\\
		=&\operatorname{Tr}_{q}^{q^2}((1+\alpha)x)+\gamma\operatorname{Tr}_{q}^{q^2}(x)+\gamma(\operatorname{Tr}_{q}^{q^2}(x))^3.\label{3d}
	\end{align}
	and
	\begin{align}
		\operatorname{Tr}_{q}^{q^2}[\delta (T^{-1}\circ f_3^{-1}(x))]=&\operatorname{Tr}_{q}^{q^2}[\delta (\beta\delta^{q}+\beta^{q}\delta)^{-1}\{\delta^{q}\operatorname{Tr}_{q}^{q^2}((1+\alpha)x)\notag\\&+\delta^{q}\gamma\operatorname{Tr}_{q}^{q^2}(x)+\delta^{q}\gamma(\operatorname{Tr}_{q}^{q^2}(x))^3+\beta^{q}\operatorname{Tr}_{q}^{q^2}(x)\}]\notag\\
		=&\operatorname{Tr}_{q}^{q^2}[\delta\delta^{q}(\beta\delta^{q}+\beta^{q}\delta)^{-1}\operatorname{Tr}_{q}^{q^2}((1+\alpha)x)\notag\\&+\delta\delta^{q}\gamma(\beta\delta^{q}+\beta^{q}\delta)^{-1}\operatorname{Tr}_{q}^{q^2}(x)\notag\\&+\delta\delta^{q}\gamma(\beta\delta^{q}+\beta^{q}\delta)^{-1}(\operatorname{Tr}_{q}^{q^2}(x))^3\notag\\&+\delta\beta^{q}(\beta\delta^{q}+\beta^{q}\delta)^{-1}\operatorname{Tr}_{q}^{q^2}(x)]\notag\\
		=&\operatorname{Tr}_{q}^{q^2}(x).\label{3e}
	\end{align}
	Now using equations \ref{3d} and \ref{3e} in equation \ref{3c}, we get
	\begin{align*}
		f_3^{-1}(x)=&\operatorname{Tr}_{q}^{q^2}((1+\alpha)x)+\gamma\operatorname{Tr}_{q}^{q^2}(x)+\gamma(\operatorname{Tr}_{q}^{q^2}(x))^3+\alpha\operatorname{Tr}_{q}^{q^2}(x)\\
		=&x+\gamma\operatorname{Tr}_{q}^{q^2}(x)+\gamma(\operatorname{Tr}_{q}^{q^2}(x))^3.
	\end{align*}
	
	Next we assume $\operatorname{Tr}(\gamma)=1$. Clearly, $\gamma\in \mathbb{F}_{q^2}\setminus\mathbb{F}_q$ and 
	$\{1,\gamma\}$ is a basis of $\mathbb{F}_{q^2}$ over $\mathbb{F}_{q}$.
	Let $U(x)=\operatorname{Tr}_q^{q^2}{(\beta x)}+\gamma\operatorname{Tr}_q^{q^2}(\delta x)$. By lemma \ref{lem2}, we have
	$U^{-1}(x)=(\beta\delta^{q}+\delta\beta^{q})^{-1}\{\delta^{q}\operatorname{Tr}_{q}^{q^2}((1+\gamma)x)+\beta^{q}\operatorname{Tr}_{q}^{q^2}(x)\}.$
	Using equation \ref{3a}, we can easily see that
  \begin{align}
  	f_3(U(x))
  	=&\operatorname{Tr}_q^{q^2}(\beta x)+\gamma(\operatorname{Tr}_q^{q^2}(\delta x))^3.\label{3f}
  \end{align}
	Next we consider the equation $f_3\circ U(x)=U(y)=\operatorname{Tr}_{q}^{q^2}(\beta y)+\gamma\operatorname{Tr}_{q}^{q^2}(\delta y)$. This gives the following system of equations	
	\begin{align*}
		\operatorname{Tr}_q^{q^2}(\beta x)=\operatorname{Tr}_q^{q^2}(\beta y)\\
		\operatorname{Tr}_q^{q^2}(\delta x)^3=\operatorname{Tr}_q^{q^2}(\delta y).
	\end{align*}
	Let $t$ be a positive integer satisfying $3t\equiv 1\mod{(2^m-1)}$.
	Then we have $\operatorname{Tr}_q^{q^2}{\delta x}=(\operatorname{Tr}_q^{q^2}(\delta y))^t$. Equivalently, we have the following system of equation
	$$\begin{bmatrix}
		\beta&\beta^{q}\\
		\delta&\delta^{q}
	\end{bmatrix}
	\begin{bmatrix}
		x\\x^{q}
	\end{bmatrix}=
	\begin{bmatrix}
		\operatorname{Tr}_{q}^{q^2}(\beta y)\\
		(\operatorname{Tr}_{q}^{q^2}(\delta y))^t
	\end{bmatrix}.$$
	Applying elementary row operation $R_1\rightarrow\delta^{q}R_1$ and $R_2\rightarrow\beta^{q}R_2$, we get
	$$\begin{bmatrix}
		\beta\delta^{q}&\beta^{q}\delta^{q}\\
		\beta^{q}\delta&\beta^{q}\delta^{q}
	\end{bmatrix}
	\begin{bmatrix}
		x\\x^{q}
	\end{bmatrix}=
	\begin{bmatrix}
		\delta^{q}\operatorname{Tr}_{q}^{q^2}(\beta y)\\
		\beta^{q}(\operatorname{Tr}_{q}^{q^2}(\delta y))^t
	\end{bmatrix}.$$
	Again applying elementary row operation $R_1\rightarrow R_1+R_2$, we get
	\begin{align}\begin{bmatrix}
			\beta\delta^{q}+\beta^{q}\delta&0\\
			\beta^{q}\delta&\beta^{q}\delta^{q}
		\end{bmatrix}
		\begin{bmatrix}
			x\\x^{q}
		\end{bmatrix}=
	\begin{bmatrix}
		\delta^{q}\operatorname{Tr}_{q}^{q^2}(\beta y)+	\beta^{q}(\operatorname{Tr}_{q}^{q^2}(\delta y))^t\\
		\beta^{q}(\operatorname{Tr}_{q}^{q^2}(\delta y))^t
	\end{bmatrix}.\label{3g}\end{align}
	From the equation \ref{3g}, we get
	\begin{align*}
		x=&(\beta\delta^{q}+\beta^{q}\delta)^{-1}\{\delta^{q}\operatorname{Tr}_{q}^{q^2}(\beta y)+\beta^{q}(\operatorname{Tr}_{q}^{q^2}(\delta y))^t\}
	\end{align*}
	or\begin{align*}
		U^{-1}\circ f_3^{-1}\circ U(y)=&(\beta\delta^{q}+\beta^{q}\delta)^{-1}\{\delta^{q}\operatorname{Tr}_{q}^{q^2}(\beta y)+\beta^{q}(\operatorname{Tr}_{q}^{q^2}(\delta y))^t\}.
	\end{align*}
	For $y=U^{-1}(x)$, we get
	\begin{align*}
		U^{-1}\circ f_3^{-1}(x)=&(\beta\delta^{q}+\beta^{q}\delta)^{-1}\{\delta^{q}\operatorname{Tr}_{q}^{q^2}(\beta U^{-1}(x))+\beta^{q}(\operatorname{Tr}_{q}^{q^2}(\delta U^{-1}(x)))^t\}.
	\end{align*}
	Now using Lemma \ref{lem3}, we get
	\begin{align*}
		U^{-1}\circ f_3^{-1}(x)=&(\beta\delta^{q}+\beta^{q}\delta)^{-1}\{\delta^{q}\operatorname{Tr}_{q}^{q^2}((1+\gamma)x)+\beta^{q}(\operatorname{Tr}_{q}^{q^2}(x))^t\}.
	\end{align*}
	Taking each side composition with $U$, we get
	\begin{align*}
		f_3^{-1}(x)=&\operatorname{Tr}_{q}^{q^2}[\beta (\beta\delta^{q}+\beta^{q}\delta)^{-1}\{\delta^{q}\operatorname{Tr}_{q}^{q^2}((1+\gamma)x)+\beta^{q}(\operatorname{Tr}_{q}^{q^2}(x))^t\}]\\&+\gamma\operatorname{Tr}_{q}^{q^2}[\delta (\beta\delta^{q}+\beta^{q}\delta)^{-1}\{\delta^{q}\operatorname{Tr}_{q}^{q^2}((1+\gamma)x)+\beta^{q}(\operatorname{Tr}_{q}^{q^2}(x))^t\}]\\
		=&\operatorname{Tr}_{q}^{q^2}((1+\gamma)x)+\gamma(\operatorname{Tr}_{q}^{q^2}(x))^t.
	\end{align*}
\end{proof}

\end{document}
